% =========================================================
% Chapter 4 — Non-stationary time series
% Author : Aymen Ben Brik (aymen.benbrik@esprit.tn)
% =========================================================

\chapter{Non-stationary time series}
\label{ch:nonstationarity}

\section*{Introduction}
\addcontentsline{toc}{section}{Introduction}

Most real-world time series are \emph{not} stationary: GDP grows,
prices drift, populations expand, atmospheric CO$_2$ accumulates.
Applying the tools of Chapter~\ref{ch:stationarity} (ACF, ARMA) to a
raw non-stationary series gives misleading results — the famous
\emph{spurious regression} phenomenon. This chapter formalises
non-stationarity, introduces the algebraic operators (backshift,
difference) and the unit-root tests (ADF, KPSS) needed to detect it,
and culminates in the \emph{ARIMA} and \emph{SARIMA} families that
combine differencing with the ARMA structure of
Chapter~\ref{ch:arma}.

% =========================================================
\section{Limits of the stationarity assumption}
\label{sec:non-stat-limits}

\begin{warnbox}[Spurious regression]
\label{warn:spurious}
If $X_t$ and $Y_t$ are two \emph{independent} random walks and one
fits the linear regression $Y_t = \beta_0 + \beta_1 X_t + u_t$ by
OLS:
\begin{itemize}
  \item the standard $t$-statistic for $\widehat\beta_1$ rejects the
        null $\beta_1 = 0$ at any level with probability tending to
        \emph{one} as the sample size grows;
  \item the $R^2$ stays large (often $> 0.5$);
\end{itemize}
even though the two series are \emph{independent}. Granger \& Newbold
(1974) called this the \emph{spurious regression} effect; it is
caused by the non-stationarity of $X$ and $Y$.
\end{warnbox}

\begin{rembox}[Practical consequence]
Before applying any classical inference (regression, ARMA, $t$-test),
\emph{check stationarity}. If the series is non-stationary, transform
it first (typically by differencing).
\end{rembox}

\subsection{Trend stationary vs.\ difference stationary}

\begin{defbox}[TS and DS processes]
\label{def:ts-ds}
\begin{itemize}
  \item \emph{Trend stationary} (TS):
        $X_t = f(t) + \varepsilon_t$ where $f$ is deterministic and
        $\varepsilon_t$ is stationary. Shocks have \emph{transient}
        effects.
  \item \emph{Difference stationary} (DS):
        $\Delta X_t$ is stationary, but $X_t$ is not. The series has
        a \emph{unit root}; shocks have \emph{permanent} effects.
\end{itemize}
\end{defbox}

\begin{rembox}
TS series are made stationary by \emph{detrending}; DS series, by
\emph{differencing}. Visual inspection alone cannot reliably
distinguish the two — a slowly trending DS process looks identical to
a TS one. We use \emph{unit-root tests} (Section~\ref{sec:unit-roots}).
\end{rembox}

% =========================================================
\section{Backshift and difference operators}
\label{sec:operators}

\begin{defbox}[Backshift operator]
\label{def:backshift}
$B X_t = X_{t-1}$, and $B^k X_t = X_{t - k}$.
For polynomials $\phi(z) = 1 - \phi_1 z - \dots - \phi_p z^p$ and
$\theta(z) = 1 + \theta_1 z + \dots + \theta_q z^q$, the AR(p) and
MA(q) processes are written
$\phi(B) X_t = \varepsilon_t$ and $X_t = \theta(B) \varepsilon_t$.
\end{defbox}

\begin{propbox}[Stationarity through the AR polynomial]
\label{prop:ar-stationary}
The AR(p) process $\phi(B) X_t = \varepsilon_t$ is (weakly) stationary
if and only if every root of $\phi(z) = 0$ lies \emph{outside} the
closed unit disk.
\end{propbox}

\begin{exbox}[AR(1)]
$\phi(z) = 1 - \phi z$ has root $z = 1 / \phi$.
$\abs{1 / \phi} > 1 \Leftrightarrow \abs\phi < 1$. We recover the
familiar AR(1) stationarity condition.
\end{exbox}

\begin{defbox}[Difference operator]
\label{def:difference}
$\Delta = 1 - B$, so $\Delta X_t = X_t - X_{t-1}$.
$\Delta^d = (1 - B)^d$ is the $d$-th difference.
The seasonal difference of period $s$ is
$\Delta_s = 1 - B^s$.
\end{defbox}

\begin{exbox}[Higher-order differencing]
$\Delta^2 X_t = (1 - B)^2 X_t = X_t - 2 X_{t-1} + X_{t-2}$.
A quadratic trend $X_t = a + b\,t + c\,t^2$ is reduced to a constant
by $\Delta^2$:
$\Delta^2 X_t = 2c$, which is stationary.
\end{exbox}

\begin{defbox}[Integration order $I(d)$]
$X_t \sim I(d)$ if $\Delta^d X_t$ is stationary but
$\Delta^{d - 1} X_t$ is not. White noise is $I(0)$; a random walk is
$I(1)$; a random walk with a stochastic trend is often $I(2)$.
\end{defbox}

% =========================================================
\section{Unit-root tests}
\label{sec:unit-roots}

\subsection{Augmented Dickey--Fuller (ADF) test}

\begin{defbox}[ADF regression]
\label{def:adf}
Starting from an AR(1) $X_t = \rho X_{t-1} + \varepsilon_t$,
substitute $\rho = 1 + \gamma$ and subtract $X_{t-1}$:
\[
  \Delta X_t = \gamma\, X_{t-1} + \varepsilon_t,
\]
augmented with $p$ lags of $\Delta X_t$ to whiten the error:
\[
  \Delta X_t = \gamma\, X_{t-1}
             + \sum_{i = 1}^{p} \alpha_i\, \Delta X_{t-i}
             + \varepsilon_t.
\]
\end{defbox}

\begin{thbox}[ADF test]
\label{th:adf}
\begin{itemize}
  \item $H_0 : \gamma = 0$ — unit root, $X_t$ is non-stationary.
  \item $H_1 : \gamma < 0$ — $X_t$ is stationary.
\end{itemize}
The $t$-statistic of $\widehat\gamma$ \emph{does not} follow a
Student's $t$ under $H_0$; its asymptotic distribution is the
\emph{Dickey--Fuller distribution}, with 5\% critical value
$\approx -2.86$ for the constant-only specification (and
$\approx -3.41$ with constant and trend).
\end{thbox}

\begin{rembox}[Three flavours]
\begin{itemize}
  \item \emph{No constant, no trend} (rare): $X_t$ should fluctuate
        around zero.
  \item \emph{Constant only} (most common): $X_t$ should fluctuate
        around a non-zero level.
  \item \emph{Constant and trend}: $X_t$ has a deterministic linear
        trend.
\end{itemize}
The choice affects the critical values; the Python function
\texttt{adfuller} accepts a \texttt{regression} argument.
\end{rembox}

\subsection{KPSS test}

\begin{defbox}[KPSS test, Kwiatkowski et al.\ 1992]
\label{def:kpss}
\[
  H_0 : \text{stationarity (around a level or a trend)},
  \qquad
  H_1 : \text{unit-root non-stationarity}.
\]
The test statistic is built from the partial sums of OLS residuals
versus a long-run variance estimator. 5\% critical values: $0.463$
for level stationarity, $0.146$ for trend stationarity. Reject $H_0$
if the statistic \emph{exceeds} the critical value.
\end{defbox}

\begin{rembox}[Combining ADF and KPSS]
The two tests have \emph{opposite} null hypotheses, which clarifies
ambiguous cases:
\begin{center}
\begin{tabular}{l|cc}
                              & ADF rejects        & ADF does not reject \\
\hline
KPSS rejects                  & contradiction      & non-stationary \\
KPSS does not reject          & stationary         & inconclusive (small $T$) \\
\end{tabular}
\end{center}
The contradiction case typically arises when there is a structural
break.
\end{rembox}

\subsection{Phillips--Perron test (mention)}

\begin{rembox}
The \emph{Phillips--Perron} (PP) test is an alternative to ADF that
adjusts the test statistic for heteroscedasticity and autocorrelation
non-parametrically (instead of via lagged differences). Same null
hypothesis as ADF, similar critical values.
\end{rembox}

% =========================================================
\section{ARIMA(p, d, q)}
\label{sec:arima}

\begin{defbox}[ARIMA model]
\label{def:arima}
$X_t \sim \ARIMA{p}{d}{q}$ if its $d$-th difference is a stationary
$\ARMA{p}{q}$:
\[
  \phi(B)\, (1 - B)^d\, X_t \;=\; \theta(B)\, \varepsilon_t,
  \qquad \varepsilon_t \sim \WN(0, \sigma^2).
\]
\end{defbox}

\begin{rembox}[Practical recipe]
\begin{enumerate}
  \item Plot $X_t$; identify visual non-stationarity.
  \item Test (ADF + KPSS); difference until stationary; choose $d$.
  \item ACF/PACF of $\Delta^d X_t$ $\Rightarrow$ candidate $(p, q)$.
  \item MLE estimate of the ARMA(p, q) coefficients.
  \item Ljung--Box on residuals; if not white, increase $p$ or $q$.
  \item Forecast, integrate back.
\end{enumerate}
\end{rembox}

\begin{exbox}[ARIMA(1, 1, 1)]
\label{ex:arima-111}
$\phi(B) = 1 - \phi B$, $\theta(B) = 1 + \theta B$, $d = 1$:
\[
  (1 - \phi B)(1 - B) X_t = (1 + \theta B)\, \varepsilon_t.
\]
With $W_t = \Delta X_t$:
$W_t = \phi W_{t-1} + \varepsilon_t + \theta \varepsilon_{t-1}$,
i.e.\ $W_t \sim \ARMA{1}{1}$. Starting from a known initial value
$X_0$, integrate back: $X_t = X_0 + \sum_{s = 1}^{t} W_s$.
\end{exbox}

% =========================================================
\section{Seasonal ARIMA (SARIMA)}
\label{sec:sarima}

When the series exhibits seasonality of period $s$, classical
differencing $(1 - B)^d$ alone is not enough — we also apply
seasonal differencing $(1 - B^s)^D$ and seasonal ARMA polynomials.

\begin{defbox}[SARIMA(p, d, q)(P, D, Q)$_s$]
\label{def:sarima}
\[
  \phi(B)\, \Phi(B^s)\,
  (1 - B)^d\, (1 - B^s)^D\, X_t
  \;=\;
  \theta(B)\, \Theta(B^s)\, \varepsilon_t,
\]
where $\Phi(z) = 1 - \Phi_1 z - \dots - \Phi_P z^P$ and
$\Theta(z) = 1 + \Theta_1 z + \dots + \Theta_Q z^Q$ are the
\emph{seasonal} AR and MA polynomials.
\end{defbox}

\begin{exbox}[SARIMA(1, 1, 1)(1, 1, 1)$_{12}$]
\[
  (1 - \phi B)(1 - \Phi B^{12})(1 - B)(1 - B^{12}) X_t
  = (1 + \theta B)(1 + \Theta B^{12}) \varepsilon_t.
\]
This model is the workhorse for monthly time series with both
short-term persistence and yearly seasonality (sales, climate, energy
demand).
\end{exbox}

\begin{rembox}[Identification]
After the differencing step, look at the ACF/PACF of
$(1 - B)^d (1 - B^s)^D X_t$:
\begin{itemize}
  \item significant spikes at lags $s, 2s, 3s, \dots$ in the ACF
        $\Rightarrow$ seasonal MA component (order $Q$);
  \item spikes at $s, 2s, \dots$ in the PACF $\Rightarrow$ seasonal
        AR component (order $P$);
  \item spikes at small lags $\Rightarrow$ ordinary $p, q$.
\end{itemize}
\end{rembox}

% =========================================================
\section*{Summary}
\addcontentsline{toc}{section}{Summary}

\noindent\textbf{Key takeaways.}
\begin{itemize}
  \item Most real series are non-stationary; \emph{spurious
        regression} can fool naive applications of OLS or ARMA.
  \item Distinguish \emph{trend stationary} (detrend) from
        \emph{difference stationary} (difference) using the ADF and
        KPSS tests.
  \item The backshift $B$ and difference $\Delta$ operators give a
        compact algebra for non-stationary models; integration order
        $I(d)$ counts the number of differencings needed.
  \item \emph{ARIMA(p, d, q)} = ARMA on the $d$-th difference;
        \emph{SARIMA} adds seasonal AR, MA and differencing.
  \item In Python, \texttt{statsmodels.tsa.stattools.adfuller} and
        \texttt{kpss}, plus \texttt{SARIMAX}, cover the whole
        workflow.
\end{itemize}
